\chapter{Evaluation and Validation}
\label{ch:evaluation-validation}

\section{Introduction to Evaluation}
\label{sec:evaluation-introduction}

This chapter evaluates whether the implemented system satisfies the central hypothesis of this thesis: automated robot-program generation from workpiece geometry should reduce preparation effort while preserving feasible execution in the simulated press-brake cell. The evaluation therefore does not focus only on whether a motion sequence can be generated. It also examines how much effort is saved, how the system behaves when the workpiece geometry or position changes, how regripping and opposite-side bending affect process time, and how the system reacts when a scenario cannot be completed safely.

The evaluation is restricted to the simulation-based scope defined in Chapter~\ref{ch:introduction}. Consequently, the chapter validates robot-program generation, simulated execution feasibility, planning behaviour, regrip handling, and safe stopping behaviour. It does not claim validation of physical bending force, material springback, final bend-angle accuracy, gripper slip on real sheet metal, or collision-free operation in a commissioned industrial cell. This distinction is important because the implemented system can be judged rigorously within its intended scope without overstating its readiness for physical production.

The analysis is organised around two complementary evaluation layers. First, manual and automated workflows are compared using repeated workflow trials in order to quantify preparation time, execution time, and total process time. Second, automated feasibility trials are evaluated scenario by scenario to assess success rate, failure rate, planning time, and controlled-stop behaviour. Separating these two layers prevents the time-comparison results from being mixed with automated-only feasibility results and keeps the interpretation of each metric consistent.

\section{Evaluation Methodology}
\label{sec:evaluation-methodology}

\subsection{Measurement Principles}
\label{subsec:measurement-principles}

All reported values are evaluated as trial-based measurements. Scenario-wise averages are interpreted first because the difficulty of the task changes with the selected bend side, workpiece geometry, workpiece position, and regrip requirement. Overall averages are then used only as compact summary indicators. This avoids treating the complete trial set as if all cases had the same technical difficulty.

The arithmetic mean is used to obtain the representative value for each scenario. For a measured quantity \(x_i\) in trial \(i\), the average value \(\bar{x}\) is calculated as:

\begin{equation}
\bar{x}=\frac{1}{N}\sum_{i=1}^{N}x_i .
\label{eq:evaluation-average}
\end{equation}

In this equation, \(N\) denotes the number of trials in the evaluated scenario and \(x_i\) denotes the measured value in one trial. This equation is used for preparation time, execution time, total process time, planning time, and regrip count. It is suitable here because the objective is not to model a probability distribution, but to compare repeated workflow observations under defined scenario conditions.

\phantomsection\label{subsec:preparation-time-definition}
\noindent\textbf{Preparation time.}
Preparation time measures the effort required before a workflow is ready for simulated execution. In the manual workflow, measurement starts when the operator begins adapting the robot task through waypoint, function-block, or sequence changes. It ends when the manual workflow is execution-ready. In the automated workflow, measurement starts when the operator begins entering or updating the workpiece geometry, pose, and bend-side input. It ends when the system has generated an execution-ready workflow. The measurement includes input selection, configuration updates, manual waypoint editing where applicable, function-block adaptation where applicable, and preparation of regrip or flip logic where applicable. It excludes simulated robot motion, physical bending, post-processing, documentation, and figure preparation. The comparative evaluation uses 48 manual workflow trials and 48 automated workflow trials across direct bending, opposite-side bending, geometry variation, and position-change conditions.

\phantomsection\label{subsec:execution-time-definition}
\noindent\textbf{Execution time.}
Execution time measures the simulated robot runtime after the workflow is ready. Measurement starts when the prepared or generated workflow begins execution and ends when the simulated sequence reaches its final state, completes unloading, or stops in a controlled way. It includes robot approach, gripping, transfer motion, bending-position motion, regrip motion, flip motion, retreat, and unloading where applicable. It excludes preparation, manual programming, operator input before execution, physical material deformation, and post-processing. Execution time is therefore interpreted as the simulated runtime of the prepared process, not as the full engineering effort.

\phantomsection\label{subsec:total-process-time-definition}
\noindent\textbf{Total process time.}
Total process time combines preparation and execution because a workflow that executes quickly but requires long manual preparation may still be inefficient in practice. For each trial, total process time is calculated as:

\begin{equation}
t_{\mathrm{total}}=t_{\mathrm{prep}}+t_{\mathrm{exec}} .
\label{eq:total-process-time}
\end{equation}

Here, \(t_{\mathrm{total}}\) is the total process time, \(t_{\mathrm{prep}}\) is the preparation time, and \(t_{\mathrm{exec}}\) is the simulated execution time. This equation is used to compare the complete effort from the start of task preparation to the end of simulated execution. It is the most important time metric in this chapter because the thesis objective is to reduce programming and preparation effort, not only robot movement time.

\phantomsection\label{subsec:planning-time-definition}
\noindent\textbf{Planning time.}
Planning time is measured only for automated feasibility trials. It starts when the automated planning procedure is triggered after valid inputs are available and ends when the system has either produced a validated sequence or rejected the case. It includes orientation evaluation, grip selection, reachability decision-making, regrip decision-making, flip decision-making where required, and task-sequence generation. It excludes operator input, manual editing, robot motion, communication after dispatch, and physical bending. Ten repeated runs are evaluated for each automated feasibility scenario. Planning time is included as a secondary metric because it remains below one second in the evaluated scenarios; the practical workflow duration is dominated by preparation and execution time.

\phantomsection\label{subsec:success-failure-definition}
\noindent\textbf{Success rate, failure rate, and controlled stop.}
Automated feasibility is evaluated using ten repeated runs per scenario. A trial is classified as successful when the intended scenario objective is completed without manual waypoint correction, unresolved feasibility violation, or uncontrolled interruption. In an opposite-side workflow, success also requires that the system inserts the required flip and regrip behaviour correctly.

A trial is classified as failed when the intended execution objective is not completed because of unresolved workspace, fixture, reachability, or feasibility constraints. The success rate \(R_{\mathrm{success}}\) is calculated as:

\begin{equation}
R_{\mathrm{success}}=\frac{N_{\mathrm{success}}}{N_{\mathrm{total}}}\cdot100 .
\label{eq:success-rate}
\end{equation}

In this equation, \(N_{\mathrm{success}}\) is the number of successful trials and \(N_{\mathrm{total}}\) is the number of evaluated trials in the scenario. The failure rate is calculated in the same way using the number of failed trials.

A controlled stop is treated separately in the interpretation. It occurs when the system detects that safe continuation is not possible and prevents the workflow from proceeding. This is a system safety success because unsafe continuation is avoided, but it is also an execution failure because the intended process is not completed. This distinction is used in the requirement verification because safe stopping is evidence for validation logic, even when the process result itself is incomplete.

\phantomsection\label{subsec:regrip-count-definition}
\noindent\textbf{Regrip count.}
Regrip count measures how many intermediate regrip operations are required after the initial pickup. The first normal grip used to pick the sheet is not counted as a regrip. Measurement starts at workflow execution start and ends at completion or controlled stop. Regrip count is averaged scenario-wise and is interpreted as an indicator of workflow complexity. A higher regrip count usually increases execution time because each regrip adds approach, release, repositioning, re-approach, gripping, and validation steps.

\section{Experimental Setup and Scenarios}
\label{sec:experimental-setup-scenarios}

All trials were conducted in the implemented virtual press-brake cell described in Chapter~\ref{ch:implementation}. The evaluation uses the UR5e model, the configured \ac{tcp}, the virtual press brake, the two-cup vacuum gripping concept, and the same general rectangular sheet-metal assumptions used during implementation. The process is evaluated as a simulated right-angle bending workflow; physical forming effects such as springback and bending force are outside the measurement scope.

Five scenario classes are evaluated. The first class is the nominal Side~A direct-bending case. Side~A represents a primary-side bend that is accessible from the initial workpiece orientation. It is the simplest valid case because the robot can generate a direct handling and bending sequence without flipping the sheet. This scenario verifies baseline program generation, bend-side interpretation, grip feasibility, and simulated execution.

The second class is the Side~E opposite-side workflow. This is technically more demanding because the required bending face is not accessible to the press-brake tool in the initial orientation. Direct bending is therefore not possible: the tool would approach the wrong side of the sheet, and the resulting pose would not correspond to the required bend direction. To make the opposite face accessible, the sheet must be transferred to a flip configuration, released, re-approached, and gripped again from a suitable side. The regrip is required because the original grasp cannot simultaneously preserve reachability, maintain a valid sheet orientation, and allow the subsequent opposite-side bend. As a result, Side~E workflows contain additional transfer, release, re-approach, flip, regrip, and validation steps. These steps are expected to increase execution time compared with Side~A, but they are necessary for feasibility.

The third class covers workpiece geometry variation. In these trials, the sheet dimensions and thickness conditions vary within the rectangular workpiece scope. This scenario tests whether the system adapts by updating input parameters rather than requiring manual reteaching of the robot path. Geometry variation is important because changes in size and thickness affect grip-point feasibility, motion distance, and reachability.

The fourth class covers workpiece position changes. Two initial placement conditions are evaluated and referred to as Plane~A and Plane~B in the figures. Plane~B represents a more constrained placement condition and is therefore expected to increase process difficulty. This scenario tests whether the automated system can absorb a placement change through parameterised planning, while the manual workflow requires greater reprogramming effort.

The fifth class covers failure-prone conditions. These trials intentionally test boundary behaviour where the system may not complete the workflow. Such cases are necessary because a strong evaluation must show not only successful operation, but also how the system behaves when the available workspace or fixture condition prevents feasible continuation.

\section{Data Analysis and Graph Interpretation}
\label{sec:data-analysis-graph-interpretation}

\subsection{Overall Time Performance}
\label{subsec:overall-time-performance}
\label{sec:comparison-with-manual-programming}

The main performance question is whether automation reduces the complete process effort. For this reason, the first comparison uses total process time rather than execution time alone. Across the 48 manual workflow trials, the average preparation time is \SI{301.0}{\second}, the average execution time is \SI{40.5}{\second}, and the average total process time is \SI{341.5}{\second}. Across the 48 automated workflow trials, the average preparation time is \SI{11.1}{\second}, the average execution time is \SI{69.7}{\second}, and the average total process time is \SI{80.7}{\second}.

These values show an important engineering trade-off. The automated workflow does not minimise every simulated robot movement. In fact, its average execution time is higher because the generated process includes structured approach motions, validation-driven transitions, and additional handling steps in complex scenarios. However, the preparation time decreases from \SI{301.0}{\second} to \SI{11.1}{\second} on average. As shown in Figure~\ref{fig:ch6-total-process-time}, this preparation-time reduction dominates the complete process result and reduces the average total process time from \SI{341.5}{\second} to \SI{80.7}{\second}.

\begin{figure}[H]
    \centering
    \includegraphics[width=\textwidth,height=0.72\textheight,keepaspectratio]{01_total_process_time_with_preparation.png}
    \caption{Comparison of mean total process time including preparation for manual and automated workflows. The figure separates preparation time and simulated execution time, showing that the automated workflow reduces the total process time mainly by reducing preparation effort.}
    \label{fig:ch6-total-process-time}
\end{figure}

Figure~\ref{fig:ch6-total-process-time} supports the main hypothesis of this thesis. The relevant improvement is not that the robot always moves faster in the automated workflow; rather, the system avoids the large manual preparation effort required for waypoint and sequence adaptation. The average total-time reduction is approximately \SI{76.4}{\percent}, calculated from the manual and automated total process times. The average preparation-time reduction is even larger, approximately \SI{96.3}{\percent}. This indicates that the primary benefit of the proposed system lies in reducing engineering effort before execution.

\subsection{Effect of Workpiece Position}
\label{subsec:effect-workpiece-position}

The position-change scenario evaluates whether the workflow remains efficient when the initial placement of the workpiece changes. The two placement conditions differ in reachability difficulty. Plane~A represents the easier placement condition, while Plane~B introduces stronger workspace and fixture constraints. If a workflow is highly sensitive to initial placement, the process time should increase significantly when moving from Plane~A to Plane~B.

The measured averages show that the automated workflow is less sensitive to this position change. In the automated case, the average total process time increases from \SI{74.5}{\second} for Plane~A to \SI{87.0}{\second} for Plane~B. This is an increase of approximately \SI{16.8}{\percent}. In the manual workflow, the average total process time increases from \SI{180.5}{\second} to \SI{502.4}{\second}, corresponding to an increase of approximately \SI{178.3}{\percent}. Figure~\ref{fig:ch6-position-effect} shows this difference directly.

\begin{figure}[H]
    \centering
    \includegraphics[width=\textwidth,height=0.72\textheight,keepaspectratio]{03_effect_workpiece_position_total_process_time.png}
    \caption{Effect of initial workpiece position on mean total process time for manual and automated workflows. Plane~A and Plane~B denote two placement conditions with different reachability constraints.}
    \label{fig:ch6-position-effect}
\end{figure}

The trend in Figure~\ref{fig:ch6-position-effect} demonstrates that the manual workflow is strongly affected by position changes because a new placement requires additional human reasoning, waypoint adaptation, and sequence verification. The automated workflow is also affected, but the increase is much smaller because the changed placement is handled through parameterised planning. This result supports the adaptability objective defined in Chapter~\ref{ch:system-concept-methodology}. It also explains why total process time is the more meaningful metric: position changes influence both preparation and execution, not only robot runtime.

\subsection{Regrip and Opposite-Side Workflow Complexity}
\label{subsec:regrip-opposite-side-analysis}

Regripping is not interpreted as an error in this evaluation. It is a planned response to reachability and orientation constraints. In Side~A direct bending, the selected face is accessible from the initial orientation, so the workflow normally avoids intermediate regripping. In Side~E opposite-side bending, the sheet must be flipped and regripped because the required face is initially inaccessible. Regrip actions therefore increase workflow complexity, but they also enable tasks that would otherwise be infeasible.

The relationship between regrip count and process time is visible in the constrained placement cases. A regrip adds at least one additional sequence of approach, release, retreat, re-approach, grip confirmation, and continuation motion. For this reason, scenarios with higher regrip counts are expected to have longer execution times. Figure~\ref{fig:ch6-regrip-count} shows the mean regrip count grouped by workflow and placement condition.

\begin{figure}[H]
    \centering
    \includegraphics[width=\textwidth,height=0.72\textheight,keepaspectratio]{05_regrip_count_by_workflow_position.png}
    \caption{Mean regrip count for manual and automated workflows under different workpiece placement conditions. The count excludes the initial pickup and represents intermediate regrips required to maintain reachability and orientation feasibility.}
    \label{fig:ch6-regrip-count}
\end{figure}

Figure~\ref{fig:ch6-regrip-count} shows that constrained placement increases the need for regripping. The automated workflow has a low mean regrip count in Plane~A and a clearly higher value in Plane~B. This is consistent with the time results in Figure~\ref{fig:ch6-position-effect}: Plane~B is more difficult, and additional regrip actions increase execution time. The conclusion is therefore not that regripping is undesirable in itself. Instead, the result shows that regripping is the mechanism by which the system preserves feasibility when direct handling is not sufficient.

The same reasoning applies to Side~E. Direct bending cannot be used because the required face is not available to the press-brake tool. The flip sequence changes the face orientation, and the regrip establishes a new grasp that is compatible with the subsequent bend. The measured increase in execution time for complex workflows is therefore technically justified by the added manipulation steps. This is an important distinction: longer execution time in Side~E does not indicate inefficient planning alone; it reflects a physically necessary handling sequence within the defined cell.

\subsection{Planning-Time Interpretation}
\label{subsec:planning-time-interpretation}

Planning time is evaluated as a secondary metric because it represents computational overhead rather than operator effort or robot runtime. The automated feasibility trials show that planning remains below one second on average in all evaluated scenario groups. The nominal case requires approximately \SI{0.42}{\second} on average, arbitrary initial orientation requires \SI{0.48}{\second}, the workspace-constrained case requires \SI{0.54}{\second}, the regrip-required case requires \SI{0.58}{\second}, and the opposite-side workflow requires approximately \SI{0.61}{\second}. The failure-prone case averages approximately \SI{0.46}{\second} because infeasible continuation is detected without completing a longer motion sequence.

These values show that computational planning time is not the limiting factor in the evaluated system. The practical benefit comes from reducing manual preparation time, while the practical cost appears mainly in additional simulated execution time for regrip and flip operations. For this reason, planning time is discussed as evidence of computational feasibility, but it is not used as the primary productivity metric.

\section{Requirement Verification}
\label{sec:requirement-verification}

The requirement verification links the measured outcomes to the functional, non-functional, and software requirements defined in Chapter~\ref{ch:system-concept-methodology}. Tables~\ref{tab:ch6-functional-verification} and~\ref{tab:ch6-nfr-sr-verification} summarise the verification logic. Requirements that are mainly implementation-based are still evaluated here by linking them to the scenario in which the implemented behaviour becomes observable.

\begin{table}[htbp]
\centering
\normalsize
\setlength{\tabcolsep}{3pt}
\renewcommand{\arraystretch}{1.12}

\begin{tabularx}{\textwidth}{L{1.6cm} L{3.6cm} Y L{3.2cm}}
\hline
\textbf{Req.} & \textbf{Verification Scenario} & \textbf{Observed Outcome} & \textbf{Evaluation Result} \\
\hline
FR-00 & Side~A, geometry variation, position change & For valid evaluated inputs, the automated workflow reaches execution-ready state without manual waypoint construction; average automated preparation is \SI{11.1}{\second}. & Verified within simulation scope. \\
\hline
FR-01 & Workspace-constrained and failure-prone cases & Infeasible or constrained cases are detected before unsafe continuation; failure-prone cases complete in 8 of 10 runs and stop safely in unresolved cases. & Verified for warning feedback, detection, and controlled response, not for universal recovery. \\
\hline
FR-02 & Side~A and Side~E & Bend-side selection changes the generated workflow: Side~A uses direct bending, while Side~E requires flip and regrip logic. & Verified. \\
\hline
FR-03 & Regrip-required and Side~E cases & Regrip-required automated trials complete in 10 of 10 runs; Side~E uses regrip as part of the flip workflow and stops safely if the remaining constraints cannot be resolved. & Verified for evaluated cases. \\
\hline
FR-04 & Arbitrary initial orientation & Automated feasibility trials complete in 10 of 10 runs with approximately \SI{0.48}{\second} planning time on average. & Verified for tested orientation variations. \\
\hline
FR-05 & Side~E opposite-side workflow & Opposite-side workflow completes in 9 of 10 runs; one run stops in a controlled way because the remaining fixture/workspace condition is not resolvable. & Partly verified; capability shown, robustness gap remains. \\
\hline
FR-06 & Side~A, Side~E, geometry variation & Successful trials indicate that grip feasibility checks provide executable grasp selections; unresolved cases are rejected rather than continued blindly. & Verified within simplified gripper and geometry assumptions. \\
\hline
FR-07 & Failure-prone conditions & Controlled stops prevent unsafe continuation when constraints cannot be resolved. & Verified as safe behaviour; execution completion is not guaranteed in boundary cases. \\
\hline
\end{tabularx}

\caption{Functional requirement verification based on evaluated scenarios and observed trial outcomes.}
\label{tab:ch6-functional-verification}
\end{table}

\begin{table}[htbp]
\centering
\normalsize
\setlength{\tabcolsep}{3pt}
\renewcommand{\arraystretch}{1.12}

\begin{tabularx}{\textwidth}{L{1.8cm} L{3.8cm} Y L{3.2cm}}
\hline
\textbf{Req.} & \textbf{Verification Scenario} & \textbf{Observed Outcome} & \textbf{Evaluation Result} \\
\hline
NFR-01 & Manual versus automated preparation & Average preparation time decreases from \SI{301.0}{\second} to \SI{11.1}{\second}; manual waypoint construction is not required in the automated workflow. & Verified. \\
\hline
NFR-02 & Geometry and position variation & The automated workflow adapts to changed geometry and placement through configuration and operator input updates without source-code modification. & Verified within the tested rectangular sheet scope. \\
\hline
NFR-03 & Failure-prone cases & Unresolved unsafe continuation results in controlled stop behaviour. & Verified for evaluated failure cases. \\
\hline
NFR-04 & Repeated automated scenario runs & Nominal, orientation, workspace-constrained, and regrip-required scenarios each produce consistent completion classifications across 10 of 10 runs. & Verified as deterministic behaviour for tested conditions; broad statistical repeatability is not claimed. \\
\hline
NFR-05 & Position-change and constrained placement & Plane~B increases process difficulty and exposes fixture/workspace sensitivity; the validation logic detects constrained continuation and supports controlled operator decisions. & Partly verified; constraint detection is shown, recovery remains limited. \\
\hline
SR-01--SR-03 & All automated execution scenarios & Generated workflows execute in the virtual runtime using the configured communication and \ac{tcp} assumptions. & Verified within the implemented virtual cell. \\
\hline
\end{tabularx}

\caption{Non-functional and software requirement verification based on evaluated scenarios and observed trial outcomes.}
\label{tab:ch6-nfr-sr-verification}
\end{table}


The verification shows that the strongest evidence is available for automated program generation, preparation-time reduction, direct bending, orientation handling, regrip insertion, and simulated execution in nominal conditions. The weaker evidence appears in opposite-side and failure-prone scenarios, where the system can identify and stop unsafe continuation but cannot always recover to a completed process. This distinction is important for an honest evaluation: safe failure handling supports the validation requirements, but it does not remove the need for improved recovery logic.

\section{Success and Failure Analysis}
\label{sec:success-failure-analysis}

The success and failure analysis focuses on the automated feasibility trials because these trials directly test whether the generated workflow completes under different scenario conditions. Each scenario consists of ten repeated runs. Figure~\ref{fig:ch6-success-failure} is introduced here because it summarises the boundary between validated operation and unresolved constraints.

\begin{figure}[H]
    \centering
    \includegraphics[width=\textwidth,height=0.72\textheight,keepaspectratio]{07_success_failure_rate_comparison.png}
    \caption{Automated success and failure rate across six feasibility scenarios. Each scenario contains ten repeated runs and is classified according to completion, failure, or controlled-stop behaviour.}
    \label{fig:ch6-success-failure}
\end{figure}

As shown in Figure~\ref{fig:ch6-success-failure}, the nominal feasible bending case, arbitrary initial orientation case, workspace-constrained configuration, and regrip-required bending case each achieve 100\% success in the ten-run evaluation. These results indicate that the implemented pipeline is stable when the scenario remains within the intended planning envelope. The regrip-required case is particularly important because it shows that regrip insertion is not only implemented, but also executable in the tested conditions.

The opposite-side workflow achieves 90\% success, with 9 of 10 runs completing successfully. The failed run is not interpreted as uncontrolled system behaviour. It represents a case where the workflow encounters a fixture or workspace constraint after partial planning and cannot continue to a valid completion. This result is significant because Side~E is more complex than Side~A: it requires flipping, regripping, revalidation, and continuation from a changed sheet orientation. The 90\% success rate therefore validates the implemented concept, while also showing that opposite-side recovery remains less robust than direct bending.

The failure-prone scenario achieves 80\% success, with 8 of 10 runs completed and 2 unresolved cases stopped. These failures define the boundary of the current implementation. In both cases, the system prevents continuation rather than forcing an infeasible motion. This is why controlled stops are interpreted as system safety success but execution failure. From an evaluation perspective, this behaviour is preferable to uncontrolled continuation, but it still shows that the current system requires stronger recovery strategies for difficult boundary cases.

\section{Gap Analysis}
\label{sec:gap-analysis}

The evaluation confirms several strengths. First, the preparation-time reduction is substantial: the average automated preparation time is \SI{11.1}{\second}, compared with \SI{301.0}{\second} for the manual workflow. This directly supports the objective of reducing manual programming effort. Second, the automated workflow remains much less sensitive to workpiece position changes. The total process time increases by approximately \SI{16.8}{\percent} from Plane~A to Plane~B in the automated workflow, while the manual workflow increases by approximately \SI{178.3}{\percent}. Third, the automated feasibility trials show full success in the nominal, orientation, workspace-constrained, and regrip-required scenario groups.

However, the evaluation also reveals clear gaps. The first gap is that physical bending quality is not validated. The chapter evaluates simulated robot-program feasibility, but it does not measure bend angle, springback, deformation, real gripper slip, or forming force. Therefore, the results support simulation-based program generation but not production-ready bending accuracy.

The second gap is the reduced robustness of the opposite-side workflow. Side~E achieves 9 successful runs out of 10, while the nominal and regrip-required cases achieve 10 out of 10. This shows that flipping and regripping are implemented successfully in most cases, but the workflow is more sensitive to fixture and workspace constraints than direct bending.

The third gap is limited recovery in failure-prone conditions. The system correctly stops unresolved cases, but an 80\% success rate means that boundary conditions still exceed the recovery capability of the current planner. The system detects the problem, which is positive for safety, but it does not always find an alternative feasible strategy.

The fourth gap concerns the manual-versus-automated comparison. The measured time reduction is strong, but the evaluation does not isolate operator learning effects, multiple operator skill levels, or long-term usability over a large workpiece family. The result therefore supports reduced workflow effort in the evaluated scenarios, but it should not be interpreted as a complete human-factors validation.

The fifth gap is collision-free industrial deployment. The virtual execution verifies generated motion under the implemented simulation assumptions, but it does not replace physical commissioning, calibrated cell modelling, or certified safety validation. This limitation is consistent with the scope of the thesis, but it remains essential for future transfer to a real cell.

\section{Improvement Plan}
\label{sec:improvement-plan}

The improvement plan follows directly from the observed gaps.

From a technical perspective, the first priority is stronger recovery for Side~E and constrained placement cases. The current workflow can detect infeasible continuation, but it should be extended with a broader regrip search and alternative flip approaches when the first strategy fails. This would address the observed 90\% success rate in opposite-side bending and improve robustness under fixture constraints.

The second technical improvement is collision-aware path validation. The implemented system validates reachability and workflow feasibility within the virtual cell, but full collision checking would improve confidence before physical deployment. This improvement should be integrated before the command-dispatch stage so that a motion is rejected or replanned before execution.

The third technical improvement is sensor-based pose correction. Position-change trials show that placement affects process time and difficulty. A perception or feedback mechanism could reduce dependence on assumed sheet pose and allow the system to correct deviations before planning. This would strengthen both NFR-05 and the practical transferability of the system.

From a validation perspective, the next step is to increase the number and diversity of trials. Each major scenario should be repeated across more workpiece dimensions, bend-side combinations, placement conditions, and regrip configurations. The current evaluation is sufficient for simulation-based thesis validation, but larger trial matrices would allow stronger claims about robustness.

Physical validation should also be added in future work. Such validation must measure bend angle accuracy, springback, gripper slip, sheet deformation, cycle repeatability, and safety behaviour on real hardware. Without these measurements, the system should continue to be described as simulation-validated rather than production-validated.

Finally, the evaluation should be made more granular. Execution time should be separated into approach, pickup, transfer, regrip, flip, bending-positioning, retreat, and unloading components. This would make it easier to identify which part of the workflow causes the largest time increase in complex scenarios. Similarly, manual preparation trials should include multiple operators to evaluate whether the time reduction remains stable across different skill levels.

\section{Summary of Findings}
\label{sec:evaluation-summary-findings}

The evaluation supports the thesis hypothesis within the defined simulation-based scope. The automated workflow reduces average preparation time from \SI{301.0}{\second} to \SI{11.1}{\second} and reduces average total process time from \SI{341.5}{\second} to \SI{80.7}{\second}. The main benefit is therefore not faster robot motion in every case, but the removal of manual waypoint and sequence preparation effort.

Scenario analysis shows that the automated system handles nominal Side~A bending, geometry variation, position changes, orientation variation, and regrip-required cases within the evaluated conditions. Position changes affect the automated workflow much less than the manual workflow, which demonstrates the value of parameterised planning. Regrip analysis further shows that additional handling steps increase execution time, but they are necessary to maintain feasibility in constrained and opposite-side workflows.

The success and failure analysis provides a balanced result. Nominal, orientation, workspace-constrained, and regrip-required automated scenarios complete successfully in all ten evaluated runs. The Side~E opposite-side workflow completes in 9 of 10 runs, while failure-prone conditions complete in 8 of 10 runs. Controlled stops are interpreted as safe system behaviour, but they remain execution failures because the process objective is not completed.

Overall, the system demonstrates a clear reduction in preparation effort and feasible simulated execution for the tested workflow classes. At the same time, the evaluation identifies remaining gaps in physical validation, collision-aware deployment, opposite-side robustness, and boundary-case recovery. These limitations do not invalidate the implemented concept; rather, they define the next development steps required before transferring the approach from simulation-based validation to industrial commissioning.
